% !TEX root = ../main.tex
\begin{comment}
=== Revision log ===
2026-05-08  R. Cappuccio / Claude  (Wave K -- Conclusions rewrite)

  Complete rewrite of the Conclusions chapter.  The previous version
  (238 lines) contained bullet-list summaries with "Lesson:" labels,
  a TRL comparison table, separate Short-Term / Long-Term future
  research bullet lists, and an "Integration Perspective" subsection.
  The rewrite condenses the chapter into continuous prose (~120 lines
  of content), organised in five short sections: summary, cross-cutting
  observations, limitations, outlook, and a closing remark.

  Removed: tab:trl_compare (TRL table), "Industrial and Research
  Implications" section, "Integration Perspective" subsection,
  all comment-fenced dead code.

  Labels preserved: ch:conclusions.
  Labels removed: tab:trl_compare.
\end{comment}

\chapter{Conclusions and Outlook}
\label{ch:conclusions}

\section{Summary}

This dissertation has presented three studies in quantum technologies,
each addressing a practical obstacle to near-term deployment: the
cryogenic overhead of superconducting processors, the characterisation
of quantum-classical hybrid learning architectures, and the design of
single-photon detectors for high-magnetic-field environments.  The
principal findings are summarised below.

The analysis of an asymmetric-SQUID transmon qubit at 300\,GHz and
4\,K (Chapter~\ref{ch:quantum_hardware}) suggests that operation in
the pulse-tube-cooler regime is, in principle, feasible.  The thermal
population at the design point remains moderate ($n_{\mathrm{th}}
\approx 2.8\,\%$), and the projected coherence times ($T_1 \approx
15\;\mu$s, $T_2 \approx 13\;\mu$s, with $T_{2,\mathrm{echo}} \approx
26\;\mu$s under Hahn-echo refocusing) appear sufficient to support
single-qubit gate fidelities above 99.8\,\% and, under an echo
cross-resonance protocol, a two-qubit CNOT fidelity of approximately
98\,\%.  These figures fall short of the surface-code threshold by a
margin that may be bridgeable through material improvements and
optimised pulse sequences, though this remains to be demonstrated
experimentally.  If the projections are borne out, the associated
reduction in cryogenic infrastructure cost could be substantial.

The hybrid quantum-classical convolutional neural network studied in
Chapter~\ref{ch:quantum_algorithms} shows, at $R=10$ independent seeds
on a binary EuroSAT subset, a $\sim 29$-percentage-point
training-efficiency advantage at the first training epoch against a
structurally isomorphic matched-capacity classical control, paired with
a $\sim 2.6$-percentage-point asymptotic-accuracy deficit against the
same control at the final epoch ($p=2.7\times 10^{-2}$, paired Wilcoxon
exact). The credit-bounded fine-tuning campaign on the IQM Emerald
54-qubit superconducting processor confirms that parameter-shift
gradients are usable on noisy quantum hardware in the small-parameter
regime: two on-device epochs recover the full simulator validation
accuracy, on a credit-bounded budget tracked by an explicitly
calibrated cost model. A complementary paired multi-seed study at
$R=10$ seeds, on a reduced 4-qubit version of the architecture under
a calibrated IBM Heron r2 noise model, further indicates that the
hybrid pipeline is noise-resilient in the studied regime: the
classification accuracy under realistic gate-level noise is
statistically indistinguishable from the noiseless reference
(paired Wilcoxon $p=0.062$), with only a transient perturbation on
the loss visible in the first three training epochs. Whether the
early-epoch training-efficiency gain constitutes a durable,
transferable advantage or is an artefact of the specific
architecture and dataset is a question that will require broader
benchmarking to settle.

The coherent many-body model of a Rydberg-atom avalanche detector
(Chapter~\ref{ch:quantum_sensing}) predicts that single-photon
amplification ($\mathcal{A} > 3$) is maintained across external
magnetic fields spanning 0--5\,T, with dark-count rates below
0.01\,Hz at 4\,K and a detection time of order $10\;\mu$s.  These characteristics, if confirmed experimentally,
would make the scheme a candidate for the readout chain of
axion-haloscope searches, where strong solenoid fields have
traditionally constrained the choice of detector technology.

\section{Cross-cutting observations}

Three recurring themes connect the otherwise independent studies.
First, \emph{in-situ tunability} plays a central role in all three:
flux-tunable sweet spots allow the transmon frequency to be adjusted
without hardware redesign; laser retuning compensates the Zeeman shifts of the Rydberg
levels, keeping the avalanche dynamics operational across the full
0--5~T field range; variational parameters
adapt the quantum circuit to the classification task at hand.
Second, in each case the relevant figure of merit is not asymptotic
supremacy but \emph{operational efficiency}---reduced cooldown time,
fewer training iterations, suppressed dark counts---a pattern that may
reflect a more general feature of near-term quantum technologies.
Third, the simulation-driven methodology adopted throughout the thesis
proved effective for narrowing the design space before committing to
fabrication or large-scale hardware access, and the quantitative
predictions it produced are, in principle, directly testable.

\section{Limitations}

The limitations of the individual studies should be stated plainly.
The transmon design has been analysed entirely through simulation;
fabrication and measurement of NbN/AlN/NbN junctions at 300\,GHz will
be needed to test whether the loss models used here---extrapolated
from lower-frequency data---remain valid in the target regime.  The
Q-CNN results were obtained on a single benchmark dataset and a
specific circuit ansatz; the extent to which the observed
training-efficiency advantage at the first epoch and the corresponding
small asymptotic-accuracy deficit generalise to other tasks,
architectures and capacity regimes is unknown.
The Rydberg-detector model assumes idealised state preparation and
readout, and does not account for Doppler broadening, stray electric
fields, or finite laser linewidth, all of which will affect
performance in a real apparatus.

\section{Outlook}

Several natural extensions of this work can be identified.  On the
hardware side, the most pressing step is the incorporation of
realistic noise spectra---$1/f$ flux noise, charge noise, and
quasiparticle dynamics---into the transmon simulations, followed by
an experimental validation campaign on test junctions.  For the hybrid
learning architecture, systematic benchmarking on additional datasets
and alternative quantum ansätze would help delineate the boundary
between tasks that benefit from the quantum layer and those that
do not.  For the Rydberg detector, experimental measurement of the
amplification factor as a function of magnetic field strength would
provide a direct test of the central prediction.

At a broader level, it may be worth noting that the three systems
studied here operate at compatible temperature scales: the transmon
at 4\,K, the Rydberg detector in the same cryogenic environment, and
the classical optimiser at room temperature.  Whether this
compatibility could eventually be exploited to build integrated
sensing-and-processing platforms is a question that lies beyond the
scope of the present work, but it does not appear to be ruled out
by any fundamental constraint.

\section{Closing remark}

The studies collected in this dissertation do not, individually or
collectively, resolve the open problems they address.  They do,
however, suggest that simulation-driven exploration of quantum
parameter spaces can identify operating regimes of practical interest
and provide the quantitative targets that experimentalists need to
decide where to invest effort.  If the transition from laboratory
demonstrations to deployable quantum systems is to proceed
efficiently, this kind of predictive groundwork---executed with care
and assessed without overstatement---may prove to be a useful
contribution.
